\scanpage{130}
\uline{Corollaire} Soit $L$ un sous corps de $\overline{\mathbb Q}$. Pour qu'on ait
$L\subset\mathbb L$, il f. et s. que $L$ soit tot\textsuperscript{t} réel, ou extension
quadratique purement imaginaire d'un corps tot. réel.

Suffisance triviale par (iii). Nécessité : Par théorie
de Galois, $L$ est une extension de degré $1$ ou $2$
\margin{\diagram{\begin{tikzcd}[ampersand replacement=\&,row sep=small,column sep=small]\mathbb K\arrow[r,dash]\arrow[d,dash]\&\mathbb K(L)\arrow[r,dash]\arrow[d,dash]\&\mathbb L\\K=\mathbb K\cap\unsure{\mathbb L|L}\arrow[r,dash]\arrow[d,dash]\&L\&\\\mathbb Q\&\&\end{tikzcd}}}
du corps tot. réel $K=L\cap\mathbb K$, et il reste à
prouver que lorsque $L$ en est une extension
quadratique, cette dernière est tot\textsuperscript{t} imaginaire.
\deleted{\unsure{Or on}} Or \deleted{\unsure{ceci résulte du}} on est ramené à prouver
\deleted{\uline{Lemme} Soit $L$ une extension quadratique d'un corps tot.
réel $K$. Si $L$}
(quitte à remplacer $L$ par un conjugué, \inserted{ce qui ne dérange pas $\mathbb K$},) que $L$ est
imaginaire. Or s'il était réel, comme $\mathbb K$ l'est, $\mathbb K(L)$ le
serait, or $\mathbb K(L)=\mathbb L$, absurde.

\uline{NB} $\mathbb L$ est la plus gde extension galoisienne de $\mathbb Q$ telle
que dans son gpe de Galois $G'$, l'image de $\tau$ (conjugaison complexe) soit \uline{centrale}. Donc $\mathbb L$ contient l'extension \deleted{\unsure{g|gal}}\unsure{abélienne|galoisienne}
maximale de $\mathbb Q$, i.e. le corps des
racines de l'unité.
% archive sheet number: 129
